%%%%%%%%%%%%%%%%%%%%%%%%%%%%%%%%%%%%%%%%%%%%%%%%%
%%%%                ABSTRACT                 %%%%
%%%%%%%%%%%%%%%%%%%%%%%%%%%%%%%%%%%%%%%%%%%%%%%%%

\begin{abstract}
We address the TREC BioGen 2024 shared task, which requires retrieving
relevant PubMed abstracts for 65 biomedical topics and grounding
generated answers with sentence-level citations.
In this first phase we build and evaluate a complete information retrieval
pipeline over 4{,}194 PubMed abstracts, comparing five retrieval strategies:
BM25, LM-Jelinek-Mercer, LM-Dirichlet, dense KNN, and Reciprocal Rank
Fusion (RRF).  Hyperparameters are selected through query-field ablation,
encoder comparison, and 5-fold cross-validation on 32~train queries,
after which the held-out 33-query test set is evaluated exactly once.
Our best single model, KNN with the domain-specialised MedCPT encoder,
achieves NDCG@100\,=\,0.8077, MAP\,=\,0.6143, and
Recall@100\,=\,0.8933; fusing it with BM25 via RRF ($k$\,=\,60) pushes
NDCG@100 to 0.8245.
Selecting a domain-matched encoder trained on 23M PubMed click-through
pairs proves to be the single most impactful design choice,
yielding $\Delta$NDCG@100 = +0.0501 versus the general-purpose baseline.
Phases~2 and~3 will extend this retrieval backbone with cross-encoder
reranking, LLM-augmented answer generation, and a ReAct-style agentic
loop for multi-step biomedical synthesis.
\end{abstract}
